\section*{Capítulo \ref{ch:g12:equilibrium} --- Equilibrio químico: cociente y constante}

\begin{solution}{exo:g12:equilibrium:1}
(a) $Q = \dfrac{[\ce{NH3}]^2 (c^\circ)^2}{[\ce{N2}][\ce{H2}]^3}$;
(b) $Q = \dfrac{[\ce{CH3COO-}][\ce{H3O+}]}{[\ce{CH3COOH}]\,c^\circ}$ (el
agua, que es el \omterm{def:g6:solutions-and-solubility:solute}{disolvente}, no se incluye);
(c) $Q = [\ce{Ca^{2+}}][\ce{CO3^{2-}}]/(c^\circ)^2$ (el sólido no se incluye);
(d) $Q = \dfrac{[\ce{FeSCN^{2+}}]\,c^\circ}{[\ce{Fe^{3+}}][\ce{SCN-}]}$.
\end{solution}

\begin{solution}{exo:g12:equilibrium:2}
$\tau = 0.020/0.050 = 0.40$: no es total.
\end{solution}

\begin{solution}{exo:g12:equilibrium:3}
$Q = 2 < K$: sentido directo. $Q = 50 > K$: sentido inverso. $Q = K$: no
hay cambio.
\end{solution}

\begin{solution}{exo:g12:equilibrium:4}
La composición es constante, pero las reacciones directa e inversa
siguen produciéndose a la misma velocidad: las \omterm{def:g7:atoms-and-molecules:molecule}{moléculas} siguen
reaccionando en ambos sentidos.
\end{solution}

\begin{solution}{exo:g12:equilibrium:5}
No, $K$ no cambia (el \omterm{def:g10:reaction-progress-table:system}{estado final} es el mismo); el tiempo para alcanzar
el equilibrio es menor.
\end{solution}

\begin{solution}{exo:g12:equilibrium:6}
$\dfrac{x^2}{(2-x)^2} = 4$, así que $\dfrac{x}{2-x} = 2$ y
$x = \qty{1.33}{mol}$: \qty{0.67}{mol} de ácido y de \omterm{def:g11:functional-groups:alcohol}{alcohol},
\qty{1.33}{mol} de \omterm{def:g11:functional-groups:carboxylic-acid}{éster} y de agua (de nuevo, dos tercios).
\end{solution}

\begin{solution}{exo:g12:equilibrium:7}
$Q = (0.5 \times 0.5)/(0.5 \times 0.5) = 1 < 4$: sentido directo, se
forma más \omterm{def:g11:functional-groups:carboxylic-acid}{éster}.
\end{solution}

\begin{solution}{exo:g12:equilibrium:8}
Unos \qty{0.52}{mol} desde el lado del ácido y \qty{0.79}{mol} desde el
lado del \omterm{def:g11:functional-groups:carboxylic-acid}{éster}. Al cabo de unas \qty{100}{h}, las dos curvas están sobre
la línea discontinua: equilibrio.
\end{solution}

\begin{solution}{exo:g12:equilibrium:9}
Con $y$ mol hidrolizados: $\dfrac{y^2}{(1-y)^2} = \dfrac{1}{4}$, así que
$\dfrac{y}{1-y} = \dfrac{1}{2}$, $y = \frac{1}{3}$: $\frac{2}{3}$ mol de
\omterm{def:g11:functional-groups:carboxylic-acid}{éster} y de agua, $\frac{1}{3}$ mol de ácido y de \omterm{def:g11:functional-groups:alcohol}{alcohol}. La misma \omterm{def:g6:pure-substances-and-mixtures:mixture}{mezcla}
final que en la esterificación.
\end{solution}

\begin{solution}{exo:g12:equilibrium:10}
Un sólido es una fase pura, cuya ``concentración'' no cambia: no se
incluye en $Q$. Añadir más sólido no cambia $Q$, así que no desplaza el
equilibrio (mientras quede algo de sólido).
\end{solution}

\begin{solution}{exo:g12:equilibrium:11}
$K$ fija la razón entre el dióxido de carbono del gas y el del agua. Al
abrir, el gas que hay sobre el agua se escapa al \omterm{def:g4:air-a-mixture-of-gases:air}{aire}: el dióxido de
carbono del gas disminuye, $Q < K$, y sale más gas del agua, hasta que
casi no queda.
\end{solution}

\begin{solution}{exo:g12:equilibrium:12}
En el equilibrio: ácido y \omterm{def:g11:functional-groups:alcohol}{alcohol} $\frac{1}{3}$, \omterm{def:g11:functional-groups:carboxylic-acid}{éster} y agua
$\frac{2}{3}$. Retirar la mitad del agua deja $\frac{1}{3}$:
$Q = \dfrac{(2/3)(1/3)}{(1/3)(1/3)} = 2 < 4$, sentido directo. Si
reaccionan $y$ más: $(2/3 + y)(1/3 + y) = 4(1/3 - y)^2$, es decir,
$27y^2 - 33y + 2 = 0$, $y = 0.064$: el \omterm{def:g11:functional-groups:carboxylic-acid}{éster} sube a \qty{0.73}{mol}.
\end{solution}

\begin{solution}{exo:g12:equilibrium:13}
El cociente de la reacción inversa tiene el numerador y el denominador
intercambiados: $Q' = 1/Q$, así que en el equilibrio $K' = 1/K$.
Hidrólisis: $1/4 = 0.25$.
\end{solution}

\begin{solution}{exo:g12:equilibrium:14}
$Q = (2 \times 2)/(1 \times 1) = 4 = K$: la \omterm{def:g6:pure-substances-and-mixtures:mixture}{mezcla} ya está en
equilibrio; su composición no cambia.
\end{solution}

\begin{solution}{exo:g12:equilibrium:15}
$n - 0.95 = 0.9025 / (4 \times 0.05) = 4.51$, así que $n = \qty{5.5}{mol}$
de \omterm{def:g11:functional-groups:alcohol}{alcohol} por \omterm{def:g10:the-mole:amount}{mol} de ácido.
\end{solution}

\begin{solution}{pb:g12:equilibrium:1}
\textbf{1.} \ce{CH3-COOH + CH3-CH2-OH <=> CH3-COO-CH2-CH3 + H2O}.

\textbf{2.} No es total: se detiene en un equilibrio en el que ambos
sentidos siguen produciéndose.

\textbf{3.} $Q = \dfrac{n_{\text{éster}}\, n_{\text{agua}}}{n_{\text{ácido}}\, n_{\text{alcohol}}}$.

\textbf{4.} $x_{\max} = \qty{1}{mol}$; $x_f = \frac{2}{3}$ mol.

\textbf{5.} $\tau = 0.67$.

\textbf{6.} Ácido y \omterm{def:g11:functional-groups:alcohol}{alcohol}, $\frac{1}{3}$ mol cada uno; \omterm{def:g11:functional-groups:carboxylic-acid}{éster} y agua,
$\frac{2}{3}$ mol cada uno.

\textbf{7.} $K = \dfrac{(2/3)^2}{(1/3)^2} = 4$.

\textbf{8.} Un tercio hidrolizado: \omterm{def:g11:functional-groups:carboxylic-acid}{éster} y agua $\frac{2}{3}$, ácido y
\omterm{def:g11:functional-groups:alcohol}{alcohol} $\frac{1}{3}$: la misma \omterm{def:g6:pure-substances-and-mixtures:mixture}{mezcla}, $Q_{eq} = 4$.

\textbf{9.} Es la misma reacción a la misma temperatura, y $K$ solo
depende de la temperatura.

\textbf{10.} Ácido $1 - x$; \omterm{def:g11:functional-groups:alcohol}{alcohol} $3 - x$; \omterm{def:g11:functional-groups:carboxylic-acid}{éster} $x$; agua $x$.

\textbf{11.} $x^2 = 4(1 - x)(3 - x) = 4(3 - 4x + x^2)$, así que
$3x^2 - 16x + 12 = 0$.

\textbf{12.} $x = (16 \pm \sqrt{256 - 144})/6 = (16 \pm 10.58)/6$: 0.903
o 4.43; solo $x_f = \qty{0.903}{mol}$ está entre 0 y 1.

\textbf{13.} \qty{90}{\percent}.

\textbf{14.} $M = \qty{88.0}{g/mol}$: $0.903 \times 88.0 = \qty{79.5}{g}$.

\textbf{15.} Cada retirada de agua reduce el numerador de $Q$, que baja
por debajo de $K$: el sistema vuelve a avanzar en sentido directo.

\textbf{16.} La reacción continuaría hasta agotar el ácido:
\qty{100}{\percent}.

\textbf{17.} El etanol es más barato, y su exceso se separa fácilmente
del \omterm{def:g11:functional-groups:carboxylic-acid}{éster} (y se recicla).

\textbf{18.} Solo el tiempo: una reacción que casi no desprende calor
tiene una constante que apenas cambia con la temperatura, y calentar
solo hace que alcance antes el equilibrio.

\textbf{19.} Alrededor del \qty{90}{\percent} del ácido.
\end{solution}
